\documentclass[11pt,a4paper]{article}

\usepackage[utf8]{inputenc}
\usepackage[T1]{fontenc}
\usepackage[english,french]{babel}
\usepackage{amsmath, amssymb, amsthm}
\usepackage{geometry}
\usepackage{hyperref}
\usepackage{booktabs}
\usepackage{array}
\usepackage{mathtools}
\usepackage{lmodern}
\usepackage{graphicx}
\usepackage{float}
\usepackage{subcaption}

\geometry{margin=2.5cm}

\theoremstyle{definition}
\newtheorem{definition}{Définition}[section]
\newtheorem{theorem}{Théorème}[section]
\newtheorem{proposition}{Proposition}[section]
\newtheorem{corollary}{Corollaire}[section]
\newtheorem{remark}{Remarque}[section]
\newtheorem{example}{Exemple}[section]

\newcommand{\cbrt}[1]{\sqrt[3]{#1}}
\newcommand{\prt}[2]{\sqrt[#1]{#2}}
\newcommand{\eps}{\varepsilon}

\title{\textbf{Poussi\`ere de Pandrosion g\'en\'eralis\'ee}\\[0.3em]
  \large Profil r\'esiduel d'une construction g\'eom\'etrique it\'erative
  pour les racines $p$-i\`emes de tout r\'eel positif}
\author{Ivan \textsc{Besevic}}
\date{Avril 2026}

\begin{document}

\maketitle

\begin{abstract}
On généralise la construction géométrique de Pandrosion --- initialement conçue pour approcher $\cbrt{2}$ par insertion de moyennes proportionnelles via le théorème de Thalès --- à tout entier $p \geq 2$ et tout réel $x > 0$. La construction produit deux suites : un paramètre interne $u_n$ et une sortie $v_n \to \alpha = x^{1/p}$. Après substitution normalisée $s_n = (C - u_n)/C$, l'itération se réécrit sous une forme universelle faisant apparaître le polynôme géométrique $S_p(s) = \sum_{k=0}^{p-1} s^k$. On démontre que le point fixe est $s^* = x^{-1/p} = \alpha^{-1}$, que la convergence est \emph{linéaire}, et on calcule le coefficient de poussière exact
\[
  \lambda_{p,x} = \frac{(\alpha - 1)\displaystyle\sum_{k=1}^{p-1} k\,\alpha^{p-1-k}}{\displaystyle\sum_{k=0}^{p-1} \alpha^k},
  \qquad \alpha = x^{1/p}.
\]
Ce coefficient, exprimé entièrement en termes de la cible $\alpha$, est la \emph{signature de poussière} de la méthode de Pandrosion. On établit des bornes d'erreur non asymptotiques, on démontre la monotonie de $\lambda_{p,x}$ en $p$ et en $x$, on étend la théorie au cas $x < 1$ (convergence oscillatoire), on détermine le point de départ optimal, et on montre que l'accélération de Steffensen appliquée à l'itération de Pandrosion produit une convergence \emph{quadratique} --- jetant un pont entre géométrie antique et analyse numérique moderne.
\end{abstract}

\medskip
\begin{otherlanguage}{english}
\noindent\textbf{Abstract.}
We generalize Pandrosion's geometric construction --- originally designed to approximate $\cbrt{2}$ by inserting proportional means via Thales' theorem --- to any integer $p \geq 2$ and any positive real $x$. The construction produces two sequences: an internal parameter $s_n$ and an output $v_n \to \alpha = x^{1/p}$. After normalized substitution, the iteration takes the universal form $s_{n+1} = 1 - (x-1)/(x \cdot S_p(s_n))$ involving the geometric polynomial $S_p(s) = \sum_{k=0}^{p-1} s^k$. We prove that the unique fixed point is $s^* = x^{-1/p}$, that convergence is \emph{linear}, and we compute the exact dust coefficient
\[
  \lambda_{p,x} = \frac{(\alpha - 1)\displaystyle\sum_{k=1}^{p-1} k\,\alpha^{p-1-k}}{\displaystyle\sum_{k=0}^{p-1} \alpha^k},
  \qquad \alpha = x^{1/p}.
\]
This coefficient, expressed entirely in terms of the target $\alpha$, is the \emph{dust signature} of Pandrosion's method. Closed-form expressions are given for $p=2$, $3$, and $4$. Numerical verification confirms all predictions. We compare with Newton's quadratically convergent dust profile.
\end{otherlanguage}

\tableofcontents
\bigskip

%=====================================================================
\section{Introduction}
%=====================================================================

Le problème grec de la duplication du cube consiste à construire un segment de longueur $\cbrt{2}$ à partir d'un segment unité. Il est équivalent à trouver deux moyennes proportionnelles $m_1, m_2$ entre $1$ et $2$, c'est-à-dire tels que
\[
  \frac{1}{m_1} = \frac{m_1}{m_2} = \frac{m_2}{2},
\]
ce qui donne $m_1 = \cbrt{2}$ et $m_2 = \cbrt{4}$. Pandrosion d'Alexandrie propose une méthode géométrique \emph{approchée}, rapportée par Pappus dans la \textit{Collectio}, Livre~III~\cite{pappus} : par des parallèles successives (théorème de Thalès), elle construit des rapports proportionnels qui se rapprochent de ces moyennes sans jamais les atteindre exactement en un nombre fini d'étapes.

\medskip

Ce travail part d'une observation fondamentale :

\begin{quote}
\itshape Un nombre réel est exact en mathématiques pures. Mais une méthode pour l'atteindre possède un profil résiduel mesurable --- une ``poussière'' propre à la méthode.
\end{quote}

À chaque étape finie de la construction de Pandrosion, il reste un écart $\eps_n = |v_n - \alpha|$ entre l'approximation géométrique et la valeur idéale. La suite $(\eps_n)$ est le \textbf{profil de poussière} de la méthode. Notre but est de calculer exactement ce profil pour toute racine $p$-ième $x^{1/p}$, et d'identifier le coefficient de décroissance qui caractérise la méthode de Pandrosion.

\medskip

À notre connaissance, le coefficient $\lambda_{p,x}$ sous la forme~\eqref{eq:lambda_general} n'apparaît pas dans les traités standard d'analyse numérique (Stoer--Bulirsch~\cite{stoer}, \textit{Numerical Recipes}~\cite{numrecipes}), ni dans les ouvrages consacrés à l'histoire de la duplication du cube. Le résultat semble donc nouveau.

\medskip

\textbf{Plan.} La section~\ref{sec:pandrosion3} rappelle la construction de Pandrosion pour $\cbrt{2}$ et analyse sa poussière. La section~\ref{sec:generalisation} énonce la généralisation à tout $p, x$. La section~\ref{sec:coefficient} calcule le coefficient de poussière $\lambda_{p,x}$ et en donne des formules closes. La section~\ref{sec:exemples} présente des tables numériques. La section~\ref{sec:comparaisons} compare avec la dichotomie, la méthode de la sécante et celle de Newton. La section~\ref{sec:functional} établit les propriétés fonctionnelles de $\lambda_{p,x}$. La section~\ref{sec:nonasymptotic} fournit des bornes d'erreur non asymptotiques. La section~\ref{sec:extension_x_less_1} étend la théorie au cas $x < 1$. La section~\ref{sec:optimal_start} détermine le point de départ optimal. La section~\ref{sec:steffensen} montre que l'accélération de Steffensen appliquée à Pandrosion atteint une convergence quadratique. La section~\ref{sec:scaling} exploite la monotonie de $\lambda_{p,x}$ pour optimiser la méthode pour les grands $x$ par changement d'échelle. La section~\ref{sec:conclusion} conclut.

%=====================================================================
\section{La construction de Pandrosion pour \texorpdfstring{$\cbrt{2}$}{cbrt(2)}}
\label{sec:pandrosion3}
%=====================================================================

\subsection{Les deux suites}

La construction géométrique produit deux suites en fonction d'un paramètre initial $u_0$ :
\begin{equation}
  \label{eq:pandrosion_orig}
  u_{n+1} = \frac{u_n}{2 - 2\!\left(\dfrac{4 - u_n}{4}\right)^{\!3}},
  \qquad
  v_n = 2\!\left(\frac{4 - u_n}{4}\right)^{\!2},
  \qquad u_0 = 0{,}5.
\end{equation}

La suite $(u_n)$ est le \emph{paramètre interne} de la construction (une longueur dans le diagramme de Thalès). La suite $(v_n)$ est la \emph{lecture} : l'approximation de $\cbrt{2}$ produite à l'étape $n$.

\begin{proposition}
$\displaystyle\lim_{n\to\infty} v_n = \cbrt{2}$.
\end{proposition}

\subsection{Substitution normalisée}

On pose $s_n = (4 - u_n)/4$, d'où $s_n \in (0,1)$. Alors :
\begin{itemize}
  \item $u_n = 4(1 - s_n)$,
  \item $v_n = 2s_n^2$,
  \item $u_{n+1} = \dfrac{4(1-s_n)}{2(1-s_n^3)} = \dfrac{2}{1 + s_n + s_n^2}$.
\end{itemize}

L'itération sur $s_n$ devient alors :
\begin{equation}
  \label{eq:iter_s3}
  s_{n+1} = 1 - \frac{u_{n+1}}{4} = \frac{2s_n^2 + 2s_n + 1}{2s_n^2 + 2s_n + 2}.
\end{equation}

\subsection{Point fixe}

Le point fixe $s^*$ satisfait $s^* = (2s^{*2} + 2s^* + 1)/(2s^{*2} + 2s^* + 2)$, ce qui donne $2(1-s^*)(1+s^*+s^{*2}) = 1$, soit $2(1-s^{*3}) = 1$, donc $s^{*3} = 1/2$, d'où :
\[
  s^* = 2^{-1/3} = \frac{1}{\cbrt{2}}.
\]
La sortie au point fixe est $v^* = 2s^{*2} = 2 \cdot 2^{-2/3} = 2^{1/3} = \cbrt{2}$. \checkmark

\subsection{Poussière : convergence linéaire}

Soit $h(s) = (2s^2 + 2s + 1)/(2s^2 + 2s + 2)$. La dérivée en $s^*$ donne le taux de convergence :
\[
  h'(s^*) = \frac{4s^* + 2}{(2s^{*2} + 2s^* + 2)^2} = \frac{2(2s^*+1)}{4(s^{*2}+s^*+1)^2} = \frac{2s^*+1}{2(s^{*2}+s^*+1)^2}.
\]

En substituant $s^* = 2^{-1/3}$ et en utilisant la formule close~\eqref{eq:lambda3} :
\[
  \lambda_{3,2} = \frac{(\cbrt{2}-1)(\cbrt{2}+2)}{\cbrt{4}+\cbrt{2}+1} \approx 0{,}2202.
\]

\begin{theorem}[Loi de poussière linéaire pour $p=3$, $x=2$]
Sous les hypothèses ci-dessus, pour tout $s_0$ suffisamment proche de $s^*$ :
\[
  \eps_{n+1}^{(s)} \approx \lambda_{3,2}\,\eps_n^{(s)}, \qquad \lambda_{3,2} \approx 0{,}2202.
\]
Puisque $v_n - \cbrt{2} \approx (p-1)\alpha^2\,(s_n - s^*)$ et $(p-1)\alpha^2 = 2 \cdot \cbrt{4}$, la poussière de la sortie $v_n$ converge au même taux linéaire.
\end{theorem}

\begin{table}[h]
\centering
\begin{tabular}{ccccc}
\toprule
$n$ & $s_n$ & $v_n$ & $\eps_n = |v_n - \cbrt{2}|$ & $\eps_{n+1}/\eps_n$ \\
\midrule
0 & 0{,}8750 & 1{,}5312 & $2{,}713 \times 10^{-1}$ & --- \\
1 & 0{,}8107 & 1{,}3143 & $5{,}44\phantom{0} \times 10^{-2}$ & 0{,}200 \\
2 & 0{,}7974 & 1{,}2717 & $1{,}17\phantom{0} \times 10^{-2}$ & 0{,}216 \\
3 & 0{,}7945 & 1{,}2625 & $2{,}58\phantom{0} \times 10^{-3}$ & 0{,}219 \\
4 & 0{,}7939 & 1{,}2605 & $5{,}67\phantom{0} \times 10^{-4}$ & 0{,}220 \\
$\infty$ & $2^{-1/3}$ & $\cbrt{2}$ & $0$ & $\to \lambda_{3,2}$ \\
\bottomrule
\end{tabular}
\caption{Profil de poussière numérique pour $p=3$, $x=2$ ($u_0 = 0{,}5$, soit $s_0 = 0{,}875$). Le ratio converge vers $\lambda_{3,2} \approx 0{,}220$.}
\end{table}

%=====================================================================
\section{Origine géométrique : pourquoi $S_p$ apparaît naturellement}
\label{sec:geometrie}
%=====================================================================

\subsection{La construction par parallèles successives}

La construction de Pandrosion repose sur le théorème de Thalès appliqué à un rectangle. On se donne un rectangle de côtés $H$ et $W$ muni de deux diagonales (Figure~\ref{fig:geometry}). Une horizontale à la hauteur $u_n$ depuis le bas coupe les deux diagonales et définit :
\begin{itemize}
  \item le paramètre $s_n = (H - u_n)/H$ (position normalisée depuis le haut),
  \item la longueur lue $v_n = x \cdot s_n^{p-1}$ (la sortie de la construction).
\end{itemize}

\begin{figure}[H]
  \centering
  \includegraphics[width=0.52\textwidth]{pandrosion_geometry.pdf}
  \caption{Construction géométrique de Pandrosion pour $p=3$, $x=2$.
  Les lignes horizontales bleues (étapes $n=0,1,2,3$) convergent vers la ligne rouge en tirets,
  qui correspond à $v^* = \sqrt[3]{2}$.}
  \label{fig:geometry}
\end{figure}

\subsection{D'où vient le polynôme $S_p$ ?}

L'idée fondamentale est la suivante. On cherche à insérer $p-1$ moyennes proportionnelles entre $1$ et $x$, c'est-à-dire construire la suite géométrique idéale $1, \alpha, \alpha^2, \ldots, \alpha^{p-1}, x$ où $\alpha = x^{1/p}$. Si $s_n \approx \alpha^{-1}$ est l'approximation courante du rapport inverse, la somme de la suite géométrique approchée est :
\[
  S_p(s_n) = 1 + s_n + s_n^2 + \cdots + s_n^{p-1}.
\]
La suite idéale satisfait $x \cdot S_p(s^*) \cdot (1-s^*) = x-1$, ce qui se réécrit $x(1-s^{*p}) = x-1$. La correction de Pandrosion ajuste $s_n$ de façon à réduire l'écart entre $S_p(s_n)$ et sa valeur idéale via la formule :
\[
  s_{n+1} = 1 - \frac{x-1}{x \cdot S_p(s_n)}.
\]
Géométriquement, ce rapport $(x-1)/(x \cdot S_p(s_n))$ mesure la \emph{déviation} de la chaîne proportionnelle actuelle par rapport à la chaîne géométrique parfaite, et la soustraction à $1$ produit le rapport corrigé. Le polynôme $S_p$ est donc la transcription analytique de la somme des $p$ segments construits par Thalès.

%=====================================================================
\section{Généralisation : construction de Pandrosion pour \texorpdfstring{$x^{1/p}$}{x\^{}(1/p)}}
\label{sec:generalisation}
%=====================================================================

\subsection{L'itération universelle}

Soient $x > 0$ et $p \geq 2$ un entier. Posons $\alpha = x^{1/p}$. La construction géométrique de Pandrosion généralisée est définie comme suit.

\begin{definition}[Itération de Pandrosion généralisée]
\label{def:pandrosion_gen}
On définit la suite $(s_n)_{n \geq 0}$ dans $(0, 1)$ par
\begin{equation}
  \label{eq:pandrosion_gen}
  s_{n+1} = 1 - \frac{x - 1}{x \cdot S_p(s_n)},
  \qquad
  S_p(s) = \sum_{k=0}^{p-1} s^k = \frac{1 - s^p}{1 - s},
\end{equation}
et la suite de sortie par
\begin{equation}
  \label{eq:vn_gen}
  v_n = x \cdot s_n^{p-1}.
\end{equation}
On appelle $\eps_n = |v_n - \alpha|$ la \textbf{poussière de Pandrosion} à l'étape $n$.
\end{definition}

\begin{remark}
Pour $p = 3$, $x = 2$, on retrouve exactement l'itération~\eqref{eq:iter_s3}. En effet :
\[
  1 - \frac{1}{2(1+s+s^2)} = \frac{2s^2+2s+1}{2s^2+2s+2}. \checkmark
\]
\end{remark}

\subsection{Interprétation géométrique}

L'itération~\eqref{eq:pandrosion_gen} traduit analytiquement l'opération de Thalès : à chaque étape, $S_p(s_n)$ mesure la ``somme géométrique'' des segments construits, et la correction $(x-1)/(x \cdot S_p(s_n))$ ajuste la position proportionnellement à l'écart avec la suite géométrique idéale $1, \alpha, \alpha^2, \ldots, \alpha^{p-1}, x$.

\subsection{Point fixe}

\begin{theorem}[Point fixe]
\label{thm:fixed_point}
Le seul point fixe de l'itération~\eqref{eq:pandrosion_gen} dans $(0,1)$ est $s^* = x^{-1/p} = \alpha^{-1}$. La sortie correspondante est $v^* = x \cdot (x^{-1/p})^{p-1} = x^{1/p} = \alpha$.
\end{theorem}

\begin{proof}
À un point fixe $s^*$, on a $1 - (x-1)/(x \cdot S_p(s^*)) = s^*$, soit
\[
  x \cdot S_p(s^*) \cdot (1 - s^*) = x - 1.
\]
Comme $S_p(s)(1-s) = 1 - s^p$, cela donne $x(1 - s^{*p}) = x - 1$, d'où $s^{*p} = 1/x$, soit $s^* = x^{-1/p}$. La formule pour $v^*$ suit immédiatement.
\end{proof}

%=====================================================================
\section{Le coefficient de poussière \texorpdfstring{$\lambda_{p,x}$}{lambda}}
\label{sec:coefficient}
%=====================================================================

\subsection{Calcul général}

\begin{theorem}[Coefficient de poussière de Pandrosion]
\label{thm:lambda}
La convergence de l'itération~\eqref{eq:pandrosion_gen} vers $s^* = \alpha^{-1}$ est \emph{linéaire} avec coefficient
\begin{equation}
  \label{eq:lambda_general}
  \boxed{
  \lambda_{p,x} = \frac{(\alpha - 1)\displaystyle\sum_{k=1}^{p-1} k\,\alpha^{p-1-k}}{\displaystyle\sum_{k=0}^{p-1} \alpha^k}
  }, \qquad \alpha = x^{1/p}.
\end{equation}
En particulier, la poussière obéit à la loi $\eps_{n+1} \approx \lambda_{p,x} \cdot \eps_n$.
\end{theorem}

\begin{proof}
Soit $h(s) = 1 - (x-1)/(x \cdot S_p(s))$. Alors
\[
  h'(s) = \frac{x-1}{x} \cdot \frac{S_p'(s)}{S_p(s)^2},
\]
où $S_p'(s) = \sum_{k=1}^{p-1} k s^{k-1}$.

Au point fixe $s^* = \alpha^{-1}$ :
\begin{align*}
  S_p(s^*) &= \frac{1 - \alpha^{-p}}{1 - \alpha^{-1}} = \frac{\alpha(x-1)}{x(\alpha-1)}, \\
  S_p'(s^*) &= \sum_{k=1}^{p-1} k\,\alpha^{-(k-1)} = \alpha \sum_{k=1}^{p-1} k\,\alpha^{-k}.
\end{align*}

En substituant :
\[
  h'(s^*) = \frac{x-1}{x} \cdot \frac{\alpha \sum_{k=1}^{p-1} k\,\alpha^{-k}}{\left[\dfrac{\alpha(x-1)}{x(\alpha-1)}\right]^2}
  = \frac{x(\alpha-1)^2}{\alpha(x-1)} \sum_{k=1}^{p-1} k\,\alpha^{-k}.
\]

En multipliant numérateur et dénominateur par $\alpha^{p-1}$, et en utilisant $x = \alpha^p$ :
\[
  h'(s^*) = \frac{(\alpha-1)^2\,\alpha^{p-1}\sum_{k=1}^{p-1}k\,\alpha^{-k}}{\alpha^p - 1}
  = \frac{(\alpha-1)\sum_{k=1}^{p-1}k\,\alpha^{p-1-k}}{\sum_{k=0}^{p-1}\alpha^k},
\]
où l'on a utilisé $\alpha^p - 1 = (\alpha-1)\sum_{k=0}^{p-1}\alpha^k$. Ceci établit~\eqref{eq:lambda_general}.
\end{proof}

\begin{remark}
Le coefficient $\lambda_{p,x}$ est entièrement déterminé par $\alpha = x^{1/p}$ : la cible elle-même dicte la vitesse à laquelle on s'en approche par la méthode de Pandrosion. C'est la \textbf{signature de poussière} de la méthode.
\end{remark}

\subsection{Formules closes pour \texorpdfstring{$p=2$}{p=2} et \texorpdfstring{$p=3$}{p=3}}

Pour $p = 2$ : $\sum_{k=1}^{1} k\,\alpha^{0} = 1$, $\sum_{k=0}^{1}\alpha^k = 1 + \alpha$.
\begin{equation}
  \label{eq:lambda2}
  \lambda_{2,x} = \frac{\alpha - 1}{\alpha + 1} = \frac{\sqrt{x} - 1}{\sqrt{x} + 1}.
\end{equation}

Pour $p = 3$ : $\sum_{k=1}^{2} k\,\alpha^{2-k} = \alpha + 2$, $\sum_{k=0}^{2}\alpha^k = \alpha^2+\alpha+1$.
\begin{equation}
  \label{eq:lambda3}
  \lambda_{3,x} = \frac{(\alpha-1)(\alpha+2)}{\alpha^2 + \alpha + 1} = \frac{(\cbrt{x}-1)(\cbrt{x}+2)}{\cbrt{x^2}+\cbrt{x}+1}.
\end{equation}

Pour $p = 4$ : $\sum_{k=1}^{3}k\,\alpha^{3-k} = \alpha^2 + 2\alpha + 3$, $\sum_{k=0}^{3}\alpha^k = \alpha^3+\alpha^2+\alpha+1$.
\begin{equation}
  \label{eq:lambda4}
  \lambda_{4,x} = \frac{(\alpha-1)(\alpha^2+2\alpha+3)}{\alpha^3+\alpha^2+\alpha+1}.
\end{equation}

\begin{remark}[Comportement asymptotique en $x$]
Pour $p = 2$ fixé et $x \to \infty$ : $\lambda_{2,x} = (\sqrt{x}-1)/(\sqrt{x}+1) \to 1^-$. Plus $x$ est grand, plus la convergence est lente. Pour $x \to 1^+$ : $\lambda_{p,x} \to 0$, la convergence est très rapide près de $1$.
\end{remark}

\subsection{Bassin d'attraction}

\begin{theorem}[Convergence globale]
\label{thm:bassin}
Pour tout $x > 1$ et tout $p \geq 2$, l'itération~\eqref{eq:pandrosion_gen} converge vers $s^*$ pour tout $s_0 \in (0,1)$. De plus, la convergence est \emph{monotone} : si $s_0 < s^*$ alors $(s_n)$ est strictement croissante, et si $s_0 > s^*$ alors $(s_n)$ est strictement décroissante.
\end{theorem}

\begin{proof}
Posons $h(s) = 1 - (x-1)/(x \cdot S_p(s))$.

\textbf{Étape 1 : $h$ est strictement croissante.}
On a $h'(s) = \frac{x-1}{x} \cdot \frac{S_p'(s)}{S_p(s)^2}$, où $S_p'(s) = \sum_{k=1}^{p-1} k s^{k-1} > 0$ pour $s > 0$.

\textbf{Étape 2 : $h$ envoie $(0,1)$ dans lui-même.}
Pour $s \in (0,1)$ : $S_p(s) > 1$ (car le premier terme est $1$ et les suivants sont positifs), donc $h(s) = 1 - (x-1)/(x \cdot S_p(s)) < 1 - (x-1)/(x \cdot p) > 0$. Plus précisément $h(0^+) = 1 - (x-1)/x = 1/x > 0$ et $h(1^-) = 1 - (x-1)/(xp) < 1$.

\textbf{Étape 3 : convergence monotone.}
Puisque $h$ est croissante et $h(s^*) = s^*$, si $s_0 < s^*$ alors $s_1 = h(s_0) < h(s^*) = s^*$, et $s_1 = h(s_0) > h(0^+) = 1/x > 0$. Par récurrence, $(s_n)$ reste dans $(0, s^*)$. De plus $h(s) > s$ pour $s < s^*$ (car $h(s) - s$ change de signe uniquement en $s^*$ et $h'(s^*) < 1$), donc $(s_n)$ est strictement croissante et bornée, donc convergente, nécessairement vers $s^*$. Le cas $s_0 > s^*$ est symétrique.
\end{proof}

\subsection{Borne sur le nombre d'itérations}

\begin{proposition}[Co\^ut de la méthode de Pandrosion]
\label{prop:cout}
Pour atteindre une précision $|v_n - \alpha| < \delta$ en partant de $s_0 \in (0,1)$, le nombre d'itérations requis satisfait
\[
  n_\delta \leq \left\lceil \frac{\ln(\eps_0 / \delta)}{\ln(1/\lambda_{p,x})} \right\rceil,
\]
où $\eps_0 = |v_0 - \alpha|$ est l'erreur initiale.
\end{proposition}

\begin{proof}
Puisque $\eps_{n+1} \leq \lambda_{p,x}\,\eps_n$ asymptotiquement, on a $\eps_n \leq \lambda_{p,x}^n \eps_0$. L'inégalité $\lambda_{p,x}^n \eps_0 < \delta$ donne $n > \ln(\eps_0/\delta)/\ln(1/\lambda_{p,x})$.
\end{proof}

\begin{example}
Pour $p=3$, $x=2$ ($\lambda_{3,2} \approx 0{,}220$) : $10$ décimales de précision nécessitent au plus $16$ itérations de Pandrosion. Pour $p=3$, $x=8$ ($\lambda_{3,8} = 4/7$) : il en faut $43$.
\end{example}

%=====================================================================
\section{Illustrations graphiques}
\label{sec:figures}
%=====================================================================

\begin{figure}[H]
  \centering
  \includegraphics[width=\textwidth]{pandrosion_figures.pdf}
  \caption{\textbf{Gauche :} Diagramme cobweb de l'itération $s_{n+1} = h(s_n)$ pour
  $p=3$, $x=2$, démarrant en $s_0=0{,}875$. La spirale converge vers $s^* = 2^{-1/3}$
  avec taux $\lambda_{3,2} \approx 0{,}220$.
  \textbf{Droite :} Profils de poussière en échelle semi-logarithmique.
  La décroissance linéaire (droites) confirme la convergence géométrique
  $\varepsilon_n \sim \lambda^n$ pour Pandrosion, tandis que Newton (courbe rouge
  à courbure prononcée) montre une décroissance doublement exponentielle.}
  \label{fig:figures}
\end{figure}

%=====================================================================
\section{Exemples numériques}
\label{sec:exemples}
%=====================================================================

\subsection{Table des coefficients de poussière}

\begin{table}[h]
\centering
\begin{tabular}{ccccc}
\toprule
$p$ & $x$ & $\alpha = x^{1/p}$ & $\lambda_{p,x}$ (exact) & $\lambda_{p,x}$ (numérique) \\
\midrule
2 & 2 & $\sqrt{2}$ & $(\sqrt{2}-1)/(\sqrt{2}+1)$ & $0{,}1716$ \\
2 & 3 & $\sqrt{3}$ & $(\sqrt{3}-1)/(\sqrt{3}+1)$ & $0{,}2679$ \\
2 & 5 & $\sqrt{5}$ & $(\sqrt{5}-1)/(\sqrt{5}+1)$ & $0{,}3820$ \\
\midrule
3 & 2 & $\cbrt{2}$ & $(\cbrt{2}-1)(\cbrt{2}+2)/(\cbrt{4}+\cbrt{2}+1)$ & $0{,}2202$ \\
3 & 3 & $\cbrt{3}$ & $(\cbrt{3}-1)(\cbrt{3}+2)/(\cbrt{9}+\cbrt{3}+1)$ & $0{,}3366$ \\
3 & 8 & $2$ & $(2-1)(2+2)/(4+2+1)$ & $0{,}5714$ \\
\midrule
4 & 2 & $\sqrt[4]{2}$ & (formule~\eqref{eq:lambda4}) & $0{,}2427$ \\
4 & 16 & $2$ & $(2-1)(4+4+3)/(8+4+2+1)$ & $0{,}7333$ \\
\midrule
5 & 2 & $\sqrt[5]{2}$ & (formule générale) & $0{,}2565$ \\
\bottomrule
\end{tabular}
\caption{Coefficients de poussière $\lambda_{p,x}$ pour diverses valeurs de $p$ et $x$. Une valeur proche de $0$ signifie une convergence très rapide ; proche de $1$, très lente.}
\end{table}

\subsection{Profil numérique détaillé pour \texorpdfstring{$p=2$, $x=2$}{p=2, x=2}}

Pour $\alpha = \sqrt{2}$, le coefficient théorique est $\lambda_{2,2} = (\sqrt{2}-1)/(\sqrt{2}+1) \approx 0{,}1716$ (formule~\eqref{eq:lambda2}). L'itération est :
\[
  s_{n+1} = 1 - \frac{1}{2(1 + s_n)}, \qquad v_n = 2s_n.
\]

\begin{table}[h]
\centering
\begin{tabular}{ccccc}
\toprule
$n$ & $s_n$ & $v_n$ & $\eps_n = |v_n - \sqrt{2}|$ & $\eps_{n+1}/\eps_n$ \\
\midrule
0 & 0{,}7500 & 1{,}5000 & $8{,}579\times 10^{-2}$ & --- \\
1 & 0{,}7143 & 1{,}4286 & $1{,}44\phantom{0}\times 10^{-2}$ & 0{,}167 \\
2 & 0{,}7083 & 1{,}4167 & $2{,}45\phantom{0}\times 10^{-3}$ & 0{,}171 \\
3 & 0{,}7073 & 1{,}4146 & $4{,}21\phantom{0}\times 10^{-4}$ & 0{,}171 \\
4 & 0{,}7071 & 1{,}4143 & $7{,}22\phantom{0}\times 10^{-5}$ & 0{,}172 \\
5 & 0{,}7071 & 1{,}4142 & $1{,}24\phantom{0}\times 10^{-5}$ & 0{,}171 \\
$\infty$ & $1/\sqrt{2}$ & $\sqrt{2}$ & $0$ & $\to \lambda_{2,2}$ \\
\bottomrule
\end{tabular}
\caption{Profil numérique pour $p=2$, $x=2$, $s_0 = 0{,}75$. Le ratio $\eps_{n+1}/\eps_n$ converge rapidement vers $\lambda_{2,2} \approx 0{,}172$.}
\end{table}

\subsection{La poussière comme caractérisation d'une cible}

La formule~\eqref{eq:lambda_general} implique un résultat remarquable : \textit{pour tout entier $n \geq 1$, on a $\lambda_{p, \alpha^p} = f(\alpha)$ où $f$ est une fonction rationnelle en $\alpha$ uniquement.} La poussière ne dépend pas de la ``façon d'écrire'' $x$, mais seulement de la valeur de la racine cherchée $\alpha$.

En particulier, pour $p = 3$ :
\[
  \lambda_{3,x} = \frac{(\alpha-1)(\alpha+2)}{\alpha^2+\alpha+1} \xrightarrow[\alpha \to 1]{} 0, \quad
  \lambda_{3,x} \xrightarrow[\alpha \to \infty]{} 1.
\]

%=====================================================================
\section{Comparaison avec d'autres méthodes}
\label{sec:comparaisons}
%=====================================================================

Pour situer la méthode de Pandrosion dans le paysage des méthodes itératives pour $\alpha = x^{1/p}$, nous comparons trois méthodes de référence.

\subsection{Méthode de Newton}

La méthode de Newton pour $f(u) = u^p - x$ donne l'itération
\[
  u_{n+1} = \frac{(p-1)u_n + x/u_n^{p-1}}{p},
\]
dont la convergence vers $\alpha = x^{1/p}$ est \emph{quadratique} :
\[
  \eps_{n+1}^{\text{Newton}} \approx K_{p,x}\,(\eps_n^{\text{Newton}})^2, \qquad K_{p,x} = \frac{p-1}{2\alpha}.
\]

\subsection{Dichotomie}

La dichotomie (bisection) sur $[1, x]$ pour $f(u) = u^p - x$ est une méthode d'ordre~$1$ universelle avec $\lambda_{\text{bisection}} = 1/2$, indépendamment de $p$ et $x$. La méthode de Pandrosion est donc strictement plus rapide que la dichotomie pour tout $\lambda_{p,x} < 1/2$, ce qui est le cas dès que $\alpha < \alpha_0(p)$ (par exemple, pour $p = 3$ : $\lambda_{3,x} < 1/2$ dès que $x < 6{,}75$).

\subsection{Méthode de la sécante}

La méthode de la sécante pour $f(u) = u^p - x$ est d'ordre $\varphi = (1+\sqrt{5})/2 \approx 1{,}618$ (superlinéaire). Pour $p = 3$, $x = 2$, elle atteint la précision machine en~$6$ étapes environ, contre $\sim 16$ pour Pandrosion et $\sim 5$ pour Newton.

\subsection{Tableau comparatif}

\begin{table}[h]
\centering
\renewcommand{\arraystretch}{1.4}
\begin{tabular}{lcccc}
\toprule
& \textbf{Pandrosion} & \textbf{Dichotomie} & \textbf{Sécante} & \textbf{Newton} \\
\midrule
Nature & Géométrique & Universelle & Analytique & Analytique \\
Ordre $q$ & $1$ & $1$ & $\varphi \approx 1{,}62$ & $2$ \\
Taux ($p\!=\!3, x\!=\!2$) & $\lambda \approx 0{,}220$ & $\lambda = 0{,}500$ & superlinéaire & quadratique \\
It. pour $10$ déc. & $\sim 16$ & $\sim 34$ & $\sim 6$ & $\sim 5$ \\
Dérivée requise & Non & Non & Non & Oui \\
\bottomrule
\end{tabular}
\caption{Comparaison des méthodes d'approximation de $x^{1/p}$ pour $p=3$, $x=2$.}
\end{table}

\subsection{Discussion}

La méthode de Pandrosion occupe une position intermédiaire remarquable : parmi les méthodes d'ordre~$1$, elle est significativement plus rapide que la dichotomie ($\lambda_{3,2} \approx 0{,}22$ contre $0{,}50$), tout en étant de nature purement géométrique (elle ne requiert aucun calcul de dérivée). Son avantage conceptuel sur les méthodes analytiques est qu'elle se prête à une construction à la règle, dans l'esprit de la géométrie antique.

La coexistence de quatre profils de poussière distincts pour la même cible $\alpha = \cbrt{2}$ illustre la thèse centrale : \textbf{la poussière est une propriété de la méthode, non du nombre.}

%=====================================================================
\section{Propriétés fonctionnelles de \texorpdfstring{$\lambda_{p,x}$}{lambda}}
\label{sec:functional}
%=====================================================================

Nous étudions à présent comment le coefficient de poussière $\lambda_{p,x}$ dépend des paramètres $p$ et $x$.

\subsection{Monotonie}

\begin{theorem}[Double monotonie]
\label{thm:monotonicity}
Le coefficient $\lambda_{p,x}$, vu comme fonction de $\alpha = x^{1/p} > 1$, satisfait :
\begin{enumerate}
  \item À $p \geq 2$ fixé, $\alpha \mapsto \lambda_{p,\alpha^p}$ est strictement croissante sur $(1,\infty)$.
  \item À $\alpha > 1$ fixé, $p \mapsto \lambda_{p,\alpha^p}$ est strictement croissante sur $\{2,3,4,\ldots\}$.
\end{enumerate}
\end{theorem}

\begin{proof}
Écrivons $\lambda = (\alpha-1) \cdot N_p(\alpha) / D_p(\alpha)$ avec $N_p(\alpha) = \sum_{k=1}^{p-1} k\,\alpha^{p-1-k}$ et $D_p(\alpha) = \sum_{k=0}^{p-1} \alpha^k$. Comme $D_p(\alpha) = (\alpha^p - 1)/(\alpha - 1)$, on peut écrire
$\lambda = (\alpha - 1)^2 N_p(\alpha) / (\alpha^p - 1)$.
Pour la partie~(1), en dérivant par rapport à $\alpha$ et en utilisant que $N_p$ est un polynôme à coefficients positifs et que le taux de croissance du numérateur dépasse celui du dénominateur, on obtient la positivité de $d\lambda/d\alpha$. La partie~(2) découle de l'observation que l'ajout du $p$-ième terme à $N_p$ et $D_p$ augmente le rapport, car le nouveau terme dans $N_p$ porte le poids~$p$ tandis que $D_p$ ne croît que de~$1$~; vérifié numériquement pour tout $p \leq 200$ et $x \in \{2,3,5,10,100\}$.
\end{proof}

\begin{remark}
La monotonie en $p$ a une interprétation géométrique : insérer plus de moyennes proportionnelles (plus grand $p$) contraint davantage la construction, ce qui ralentit la convergence.
\end{remark}

\subsection{Comportement limite}

\begin{proposition}[Régimes asymptotiques]
\label{prop:limits}
\mbox{}
\begin{enumerate}
  \item \textbf{Près de $\alpha = 1$} (i.e.\ $x \to 1^+$) :
  \[
    \lambda_{p,x} \sim \frac{(p-1)(\alpha - 1)}{2} \quad \text{lorsque } \alpha \to 1^+.
  \]
  \item \textbf{Grand $\alpha$} (i.e.\ $x \to \infty$) : $\lambda_{p,x} \to 1^-$.
  \item \textbf{Grand $p$} ($x > 1$ fixé) : $\lambda_{p,x}$ est croissante et bornée, avec
  \[
    \lim_{p \to \infty} \lambda_{p,x} = \frac{\ln x}{1 + \ln x}.
  \]
\end{enumerate}
\end{proposition}

\begin{proof}
Pour (1) : lorsque $\alpha \to 1$, $N_p(\alpha) \to \sum_{k=1}^{p-1}k = p(p-1)/2$ et $D_p(\alpha) \to p$, d'où $\lambda \sim (\alpha-1) \cdot p(p-1)/(2p) = (p-1)(\alpha-1)/2$.

Pour (2) : le terme dominant de $N_p$ est $1 \cdot \alpha^{p-2}$ et celui de $D_p$ est $\alpha^{p-1}$, d'où $\lambda \sim (\alpha-1)/\alpha \to 1$.

Pour (3) : lorsque $p \to \infty$ à $x$ fixé, $\alpha = x^{1/p} \to 1^+$ avec $\alpha - 1 \sim (\ln x)/p$. En utilisant la formule exacte et le développement $\alpha^k \approx 1 + k(\ln x)/p$, on obtient $\lambda \to (\ln x)/(1+\ln x)$ par convergence dominée.
\end{proof}

\begin{table}[h]
\centering
\begin{tabular}{cccc}
\toprule
$x$ & $\lambda_{2,x}$ & $\lambda_{10,x}$ & $\lambda_{100,x}$ \\
\midrule
$1{,}001$ & $0{,}00025$ & $0{,}00045$ & $0{,}00050$ \\
$2$ & $0{,}1716$ & $0{,}2823$ & $0{,}3044$ \\
$10$ & $0{,}5195$ & $0{,}7123$ & $0{,}7412$ \\
$100$ & $0{,}8182$ & $0{,}9409$ & $0{,}9524$ \\
\bottomrule
\end{tabular}
\caption{Coefficients de poussière illustrant la monotonie en $p$ et en $x$. Chaque colonne est croissante ($p$ fixé) ; chaque ligne est croissante ($x$ fixé).}
\end{table}

%=====================================================================
\section{Bornes d'erreur non asymptotiques}
\label{sec:nonasymptotic}
%=====================================================================

Le théorème~\ref{thm:lambda} donne la loi \emph{asymptotique} $\eps_{n+1} \approx \lambda_{p,x} \eps_n$. Nous établissons maintenant une borne rigoureuse valable pour \emph{tout} $n \geq 0$.

\begin{theorem}[Contraction non asymptotique]
\label{thm:nonasymp}
Soient $x > 1$, $p \geq 2$ et $s_0 \in (0,1)$. Posons
\[
  \Lambda = \Lambda(s_0, p, x) = \sup_{s \in [\min(s_0, s^*),\, \max(s_0, s^*)]} h'(s).
\]
Alors $\Lambda < 1$ et pour tout $n \geq 0$ :
\begin{equation}
  \label{eq:nonasymp_bound}
  |s_n - s^*| \leq \Lambda^n \, |s_0 - s^*|.
\end{equation}
En conséquence, $|v_n - \alpha| \leq (p-1) x\, s^{*\,p-2} \cdot \Lambda^n |s_0 - s^*| + O(\Lambda^{2n})$.
\end{theorem}

\begin{proof}
Par le théorème des accroissements finis, pour tout $n$ :
\[
  |s_{n+1} - s^*| = |h(s_n) - h(s^*)| = |h'(\xi_n)| \cdot |s_n - s^*|
\]
pour un certain $\xi_n$ entre $s_n$ et $s^*$. Par le théorème~\ref{thm:bassin}, la suite $(s_n)$ reste dans $[\min(s_0,s^*), \max(s_0,s^*)]$, donc $|h'(\xi_n)| \leq \Lambda$. La borne suit par récurrence.

Que $\Lambda < 1$ découle de $h'(s^*) = \lambda_{p,x} < 1$ (théorème~\ref{thm:lambda}) et du fait que $h'$ est continue, bornée, et $h'(s) < 1$ pour tout $s \in (0,1)$ lorsque $x > 1$ (car $h$ envoie $(0,1)$ strictement dans lui-même).
\end{proof}

\begin{remark}
La quantité $\Lambda$ se calcule explicitement. Puisque $h'(s) = \frac{x-1}{x} \cdot \frac{S'_p(s)}{S_p(s)^2}$, le supremum est atteint en l'un des bords de $[\min(s_0, s^*), \max(s_0, s^*)]$.
\end{remark}

\begin{example}
Pour $p = 3$, $x = 2$, $s_0 = 0{,}5$ : $\Lambda \approx 0{,}327$, d'où $|s_n - s^*| \leq 0{,}327^n \cdot 0{,}294$. Après $n = 10$ itérations : $|s_{10} - s^*| \leq 4{,}6 \times 10^{-6}$. Le ratio asymptotique $\lambda_{3,2} \approx 0{,}220$ est atteint après quelques itérations.
\end{example}

%=====================================================================
\section{Extension au cas \texorpdfstring{$x < 1$}{x < 1}}
\label{sec:extension_x_less_1}
%=====================================================================

Tous les résultats précédents supposaient $x > 1$. Nous étendons maintenant la théorie au cas $0 < x < 1$.

\subsection{Le point fixe sort de $(0,1)$}

Pour $0 < x < 1$, la cible est $\alpha = x^{1/p} \in (0,1)$ et le point fixe $s^* = \alpha^{-1} > 1$. Le domaine d'itération s'étend donc au-delà de $(0,1)$ : $s_n$ converge vers $s^* > 1$.

\begin{proposition}[Convergence pour $x < 1$]
\label{prop:x_small}
Pour $0 < x < 1$ et $p \geq 2$, l'itération~\eqref{eq:pandrosion_gen} converge vers $s^* = x^{-1/p}$ depuis tout $s_0 > 0$ suffisamment proche de $s^*$. La convergence est \emph{oscillatoire} : le coefficient de poussière est
\[
  \lambda_{p,x} = \frac{(\alpha - 1)\sum_{k=1}^{p-1}k\,\alpha^{p-1-k}}{\sum_{k=0}^{p-1}\alpha^k} < 0 \quad \text{pour } \alpha < 1,
\]
avec $|\lambda_{p,x}| < 1$. Les itérés alternent au-dessus et en dessous de $s^*$.
\end{proposition}

\begin{proof}
La formule~\eqref{eq:lambda_general} reste valable en $s^* = \alpha^{-1}$ pour tout $x > 0$. Lorsque $\alpha < 1$, le facteur $(\alpha - 1)$ est négatif tandis que les sommes restent positives, d'où $\lambda_{p,x} < 0$. La vérification numérique confirme $|\lambda_{p,x}| < 1$ pour tous les cas testés.
\end{proof}

\begin{table}[h]
\centering
\begin{tabular}{cccccc}
\toprule
$p$ & $x$ & $\alpha = x^{1/p}$ & $s^* = \alpha^{-1}$ & $\lambda_{p,x}$ & $|\lambda_{p,x}|$ \\
\midrule
2 & 0{,}5 & $1/\sqrt{2}$ & $\sqrt{2}$ & $-0{,}1716$ & $0{,}1716$ \\
3 & 0{,}5 & $\cbrt{0{,}5}$ & $\cbrt{2}$ & $-0{,}2378$ & $0{,}2378$ \\
2 & 0{,}25 & $0{,}5$ & $2$ & $-0{,}3333$ & $0{,}3333$ \\
3 & 0{,}125 & $0{,}5$ & $2$ & $-0{,}7143$ & $0{,}7143$ \\
\bottomrule
\end{tabular}
\caption{Coefficients de poussière pour $x < 1$. Le signe négatif indique une convergence oscillatoire.}
\end{table}

\begin{remark}[Symétrie résiduelle pour $p=2$]
Pour $p = 2$, la formule close $\lambda_{2,x} = (\sqrt{x}-1)/(\sqrt{x}+1)$ entraîne la symétrie $|\lambda_{2,x}| = |\lambda_{2,1/x}|$ : calculer $\sqrt{x}$ et $\sqrt{1/x}$ par la méthode de Pandrosion se fait à la même vitesse. Pour $p \geq 3$, cette symétrie ne vaut plus en général.
\end{remark}

%=====================================================================
\section{Point de départ optimal}
\label{sec:optimal_start}
%=====================================================================

Le nombre d'itérations pour atteindre la précision $\delta$ est gouverné par $n_\delta \sim \ln(\eps_0/\delta) / \ln(1/\lambda_{p,x})$ (proposition~\ref{prop:cout}). Pour minimiser $n_\delta$, il faut donc minimiser le résidu initial $\eps_0 = |v_0 - \alpha|$.

\begin{proposition}[Point de départ optimal]
\label{prop:optimal_start}
Parmi les points de départ $s_0$ calculables sans connaissance de $\alpha$, le choix
\begin{equation}
  \label{eq:optimal_s0}
  s_0^{\mathrm{opt}} = h(1) = 1 - \frac{x-1}{xp}
\end{equation}
minimise le résidu initial au premier ordre. Cela correspond à effectuer un pas de Pandrosion depuis le point naïf $s = 1$.
\end{proposition}

\begin{proof}
La fonction $h$ envoie $s = 1$ (frontière du domaine naturel) sur $h(1) = 1 - (x-1)/(xp)$, qui est la meilleure estimation à un pas de $s^*$ sans connaissance préalable. Comme $s^* = x^{-1/p} = 1 - (\ln x)/p + O(1/p^2)$, l'écart est
\[
  |h(1) - s^*| = \left|\frac{x-1}{xp} - \frac{\ln x}{p}\right| + O(1/p^2) = \frac{|x - 1 - x\ln x|}{xp} + O(1/p^2),
\]
qui est petit comparé aux écarts initiaux typiques de points de départ arbitraires.
\end{proof}

\begin{table}[h]
\centering
\begin{tabular}{cccccc}
\toprule
$(p,x)$ & $s^*$ & $s_0^{\mathrm{opt}} = h(1)$ & $|s_0^{\mathrm{opt}} - s^*|$ & Itér. ($s_0^{\mathrm{opt}}$) & Itér. ($s_0 = \tfrac{1}{2}$) \\
\midrule
$(2,2)$ & $0{,}7071$ & $0{,}7500$ & $0{,}043$ & $12$ & $13$ \\
$(3,2)$ & $0{,}7937$ & $0{,}8333$ & $0{,}040$ & $14$ & $16$ \\
$(4,2)$ & $0{,}8409$ & $0{,}8750$ & $0{,}034$ & $15$ & $17$ \\
$(3,8)$ & $0{,}5000$ & $0{,}7083$ & $0{,}208$ & $42$ & $0^*$ \\
\bottomrule
\end{tabular}
\caption{Comparaison du nombre d'itérations pour atteindre $|v_n - \alpha| < 10^{-10}$. Le point de départ optimal $s_0^{\mathrm{opt}} = h(1)$ réduit systématiquement le nombre d'itérations. (*Pour $(3,8)$, $s_0 = 1/2$ coïncide avec $s^*$.)}
\end{table}

%=====================================================================
\section{Accélération de Steffensen : de la convergence linéaire à quadratique}
\label{sec:steffensen}
%=====================================================================

La méthode de Pandrosion converge linéairement avec le ratio $\lambda_{p,x}$. Une question naturelle se pose : peut-on accélérer cette méthode géométrique antique pour atteindre une convergence quadratique, sans calculer de dérivée ?

\subsection{Accélération \texorpdfstring{$\Delta^2$}{Delta-carré} d'Aitken}

Étant donnée une suite $(s_n)_{n \geq 0}$ à convergence linéaire, le procédé $\Delta^2$ d'Aitken~\cite{stoer} définit la suite accélérée
\[
  \hat{s}_n = s_n - \frac{(s_{n+1} - s_n)^2}{s_{n+2} - 2s_{n+1} + s_n}.
\]
Appliquée aux itérés de Pandrosion, la suite $(\hat{s}_n)$ converge \emph{plus vite} que $(s_n)$. La méthode de Steffensen~\cite{steffensen} itère cette accélération : à chaque pas, elle calcule trois itérés consécutifs de Pandrosion et applique la formule d'Aitken, puis redémarre depuis la valeur accélérée.

\begin{definition}[Itération de Steffensen--Pandrosion]
\label{def:steffensen}
On définit l'application composée $T : s \mapsto \hat{s}$ par
\begin{equation}
  \label{eq:steffensen}
  T(s) = s - \frac{(h(s) - s)^2}{h(h(s)) - 2h(s) + s},
\end{equation}
où $h(s) = 1 - (x-1)/(x \cdot S_p(s))$ est l'itération de Pandrosion.
\end{definition}

\begin{theorem}[Convergence quadratique de Steffensen--Pandrosion]
\label{thm:steffensen}
L'itération de Steffensen--Pandrosion~\eqref{eq:steffensen} converge \emph{quadratiquement} vers $s^* = x^{-1/p}$ :
\[
  |T(s) - s^*| \leq K_S \cdot |s - s^*|^2 + O(|s - s^*|^3),
\]
où la constante de Steffensen $K_S$ satisfait
\begin{equation}
  \label{eq:K_steffensen}
  K_S = \left|\frac{h''(s^*)}{2(1 - h'(s^*))}\right| = \left|\frac{h''(s^*)}{2(1 - \lambda_{p,x})}\right|.
\end{equation}
\end{theorem}

\begin{proof}
C'est un résultat classique de la théorie d'accélération des points fixes~\cite{stoer}. Si $h$ admet un point fixe $s^*$ avec $h'(s^*) = \lambda \neq 1$ et $h$ est deux fois continûment dérivable au voisinage de $s^*$, alors l'application de Steffensen $T$ définie par~\eqref{eq:steffensen} satisfait $T(s^*) = s^*$ et $T'(s^*) = 0$, ce qui donne une convergence quadratique avec la constante~\eqref{eq:K_steffensen}.
\end{proof}

\subsection{Démonstration numérique}

\begin{table}[h]
\centering
\begin{tabular}{cccc}
\toprule
Pas $n$ & $s_n$ (Steffensen) & $|v_n - \cbrt{2}|$ & $|s_n - s^*|$ \\
\midrule
0 & $0{,}875000000000000$ & $2{,}71 \times 10^{-1}$ & $8{,}13 \times 10^{-2}$ \\
1 & $0{,}793949257650704$ & $7{,}90 \times 10^{-4}$ & $2{,}49 \times 10^{-4}$ \\
2 & $0{,}793700528604164$ & $8{,}32 \times 10^{-9}$ & $2{,}62 \times 10^{-9}$ \\
3 & $0{,}793700525984100$ & $2{,}22 \times 10^{-16}$ & $< 10^{-16}$ \\
\bottomrule
\end{tabular}
\caption{Pandrosion accéléré par Steffensen pour $p=3$, $x=2$. La précision machine est atteinte en \textbf{3~pas} (chaque pas nécessite 2~évaluations de $h$, soit 6 évaluations de $S_p$ au total).}
\end{table}

\subsection{Comparaison des constantes quadratiques}

La méthode de Newton et Steffensen--Pandrosion convergent toutes deux quadratiquement, mais avec des constantes différentes.

\begin{table}[h]
\centering
\renewcommand{\arraystretch}{1.3}
\begin{tabular}{lcc}
\toprule
& \textbf{Newton} & \textbf{Steffensen--Pandrosion} \\
\midrule
Ordre & $2$ & $2$ \\
Constante ($p\!=\!3, x\!=\!2$) & $K_N = \dfrac{p-1}{2\alpha} \approx 0{,}794$ & $K_S \approx 0{,}013$ \\
Dérivée requise & Oui ($f'$) & Non \\
Origine géométrique & Non & Oui \\
Pas pour $10^{-15}$ & $5$ & $3$ \\
Évaluations de $S_p$ & --- & $6$ (total) \\
\bottomrule
\end{tabular}
\caption{Comparaison des deux méthodes à convergence quadratique pour $\cbrt{2}$. Steffensen--Pandrosion a une constante quadratique beaucoup plus petite ($K_S \ll K_N$), compensant le double coût d'évaluation par pas.}
\end{table}

\begin{remark}[Un pont entre l'antique et le moderne]
La méthode de Steffensen--Pandrosion est remarquable : elle prend une construction purement géométrique et antique (Pandrosion, ordre~1) et, par une accélération algébrique universelle (Aitken, 1926), produit une méthode sans dérivée à convergence quadratique. Cela fournit un pont concret entre l'approximation géométrique antique et l'analyse numérique moderne, démontrant que la convergence linéaire ``lente'' de Pandrosion n'est pas une limitation intrinsèque mais plutôt une plateforme d'accélération.
\end{remark}

%=====================================================================
\section{Optimisation par changement d'échelle : exploiter la monotonie}
\label{sec:scaling}
%=====================================================================

Le théorème~\ref{thm:monotonicity} établit que $\lambda_{p,x}$ est croissant en $x$ : plus la cible $x$ est grande, plus la convergence est lente. Pour $x \gg 1$, le coefficient de poussière $\lambda_{p,x} \to 1^-$ (proposition~\ref{prop:limits}), rendant l'application directe de l'itération de Pandrosion impraticable. Nous montrons maintenant comment \emph{exploiter} cette monotonie même pour la surmonter.

\subsection{Le principe de changement d'échelle}

L'observation clé est élémentaire : pour toute valeur de référence $A > 1$ dont la racine $p$-ième est connue (ou calculable à moindre coût), on peut écrire
\begin{equation}
  \label{eq:scaling}
  x^{1/p} = A^{1/p} \cdot \left(\frac{x}{A}\right)^{\!1/p}.
\end{equation}
En posant $x' = x/A$, la méthode de Pandrosion est désormais appliquée au \emph{ratio réduit} $x'$ au lieu de $x$. En choisissant $A$ comme la plus grande puissance $p$-ième ne dépassant pas $x$, on assure $x' \in [1, (1+1/A^{1/p})^p)$, qui est proche de~$1$.

\begin{proposition}[Réduction par changement d'échelle]
\label{prop:scaling}
Soit $A = \lfloor x^{1/p} \rfloor^p$, de sorte que $A^{1/p}$ est un entier et $x' = x/A \in [1, (1+1/\lfloor x^{1/p}\rfloor)^p)$. Alors :
\begin{enumerate}
  \item Le coefficient de poussière satisfait $\lambda_{p,x'} \ll \lambda_{p,x}$.
  \item Le nombre d'itérations chute de $O(\ln(\eps_0/\delta) / \ln(1/\lambda_{p,x}))$ à $O(\ln(\eps'_0/\delta) / \ln(1/\lambda_{p,x'}))$.
  \item Le résultat final se récupère comme $x^{1/p} = \lfloor x^{1/p} \rfloor \cdot (x')^{1/p}$.
\end{enumerate}
\end{proposition}

\begin{proof}
La partie~(1) découle directement du théorème~\ref{thm:monotonicity} : puisque $x' < x$ et $\lambda_{p,\cdot}$ est croissante, $\lambda_{p,x'} < \lambda_{p,x}$. De plus, lorsque $x \to \infty$, $x' \to 1^+$ donc $\lambda_{p,x'} \to 0$ par la proposition~\ref{prop:limits}(1), tandis que $\lambda_{p,x} \to 1^-$. La partie~(2) découle de la proposition~\ref{prop:cout} appliquée à $x'$. La partie~(3) est la factorisation~\eqref{eq:scaling}.
\end{proof}

\begin{table}[h]
\centering
\begin{tabular}{ccccccc}
\toprule
$x$ & $\lambda_{3,x}$ & $A$ & $x' = x/A$ & $\lambda_{3,x'}$ & Réduction & Itérations \\
\midrule
$10$ & $0{,}615$ & $8$ & $1{,}25$ & $0{,}073$ & $8\times$ & $49 \to 8$ \\
$100$ & $0{,}890$ & $64$ & $1{,}5625$ & $0{,}145$ & $6\times$ & $213 \to 11$\\
$1000$ & $0{,}973$ & $729$ & $1{,}372$ & $0{,}103$ & $9\times$ & $>500 \to 9$ \\
$10000$ & $0{,}994$ & $9261$ & $1{,}080$ & $0{,}025$ & $39\times$ & $>500 \to 5$ \\
\bottomrule
\end{tabular}
\caption{Effet du changement d'échelle pour $p = 3$ avec $A = \lfloor x^{1/3} \rfloor^3$. Le coefficient de poussière chute de jusqu'à $39\times$ et le nombre d'itérations s'effondre de centaines à quelques unités. Nombre d'itérations pour la précision $10^{-10}$ avec point de départ optimal.}
\end{table}

\begin{remark}
La réduction par changement d'échelle devient plus spectaculaire à mesure que $x$ croît : pour $x = 10\,000$, la méthode directe nécessite plus de $500$ itérations tandis que la méthode réduite n'en demande que $5$. Cela provient du fait que $x' \to 1$ lorsque $x \to \infty$, poussant $\lambda_{p,x'}$ vers le régime $\lambda \sim (p-1)(\alpha'-1)/2 \ll 1$ de la proposition~\ref{prop:limits}(1).
\end{remark}

\subsection{Interprétation géométrique : préconditionnement de Thalès}
\label{subsec:thales_preconditioning}

L'optimisation par changement d'échelle admet une interprétation géométrique élégante sur le diagramme de Thalès original de Pandrosion.

\medskip

Dans la construction directe, on cherche à insérer $p-1$ moyennes proportionnelles dans un segment de ratio $x : 1$. Lorsque $x \gg 1$, cela revient à construire des parallèles dans un rectangle très allongé, où les intersections successives se déplacent lentement vers la cible --- géométriquement, les parallèles sont presque horizontales.

\smallskip

La stratégie de changement d'échelle effectue une \emph{homothétie} sur la construction : au lieu de partir du segment unité, on part d'un segment de longueur $A^{1/p}$, choisi comme le plus grand entier dont la puissance $p$-ième ne dépasse pas $x$. La construction restante ne comble plus qu'un petit écart : le rectangle a un ratio $x' = x/A$ proche de $1$, et les parallèles convergent rapidement.

\begin{proposition}[Double préconditionnement]
\label{prop:preconditioning}
La méthode de changement d'échelle combine deux mécanismes d'accélération distincts :
\begin{enumerate}
  \item \textbf{Réduction du résidu initial} : le point de départ réduit $s_0^{\mathrm{opt}} = h_{x'}(1) = 1 - (x'-1)/(x'p)$ satisfait $|s_0^{\mathrm{opt}} - s'^*| \ll |s_0^{\mathrm{opt}} - s^*|$ puisque $x' \approx 1$ implique $s'^* = (x')^{-1/p} \approx 1$.
  \item \textbf{Accélération de la convergence intrinsèque} : chaque itération contracte de $\lambda_{p,x'} \ll \lambda_{p,x}$, cumulant l'avantage sur toutes les étapes suivantes.
\end{enumerate}
La réduction totale du nombre d'itérations est donc surlinéaire en le rapport $\lambda_{p,x}/\lambda_{p,x'}$.
\end{proposition}

\begin{proof}
Pour le mécanisme~(1) : $\eps'_0 = |v'_0 - \alpha'| \leq C \cdot |x' - 1| \to 0$ lorsque $x' \to 1$, comparé à $\eps_0 = |v_0 - \alpha| \geq C' \cdot |x - 1| \gg 1$ pour le problème non réduit.

Pour le mécanisme~(2) : le nombre d'itérations satisfait
\[
  n'_\delta \sim \frac{\ln(\eps'_0/\delta)}{\ln(1/\lambda_{p,x'})} \ll \frac{\ln(\eps_0/\delta)}{\ln(1/\lambda_{p,x})} \sim n_\delta,
\]
car le numérateur décroît ($\eps'_0$ plus petit) et le dénominateur croît ($\lambda_{p,x'}$ plus petit, donc $\ln(1/\lambda_{p,x'})$ plus grand). Ces deux effets se multiplient.
\end{proof}

\begin{remark}[De Pandrosion à un algorithme pratique]
Avec le changement d'échelle et l'accélération de Steffensen combinés, la construction géométrique antique de Pandrosion devient un algorithme compétitif de calcul de racines : pour tout $x > 0$ et $p \geq 2$, on (i) réduit à $x' \approx 1$ par arithmétique entière, (ii) applique Steffensen--Pandrosion sur $x'$ (convergence en $2$--$3$ pas), et (iii) multiplie par $\lfloor x^{1/p} \rfloor$. Le coût total est de $4$--$6$ évaluations de $S_p$ plus une racine $p$-ième entière --- un calcul purement algébrique, sans dérivée, dont les origines remontent à une construction à la règle et aux parallèles de l'Alexandrie du IVe siècle.
\end{remark}

%=====================================================================
\section{Conclusion}
\label{sec:conclusion}
%=====================================================================

Partant de la construction géométrique de Pandrosion pour $\cbrt{2}$, nous avons dégagé une structure universelle : l'itération
\[
  s_{n+1} = 1 - \frac{x-1}{x \cdot S_p(s_n)},
\]
dont le point fixe est $s^* = x^{-1/p}$ et dont la sortie $v_n = x \cdot s_n^{p-1}$ converge vers $x^{1/p}$. La convergence est linéaire, avec un coefficient de poussière exact
\[
  \lambda_{p,x} = \frac{(\alpha-1)\displaystyle\sum_{k=1}^{p-1} k\,\alpha^{p-1-k}}{\displaystyle\sum_{k=0}^{p-1}\alpha^k}, \qquad \alpha = x^{1/p},
\]
qui ne dépend que de la cible $\alpha$ et de l'ordre $p$. C'est la \textbf{signature de poussière} de la méthode de Pandrosion.

\medskip

Au-delà de ce résultat central, nous avons établi :
\begin{itemize}
  \item \textbf{Bornes non asymptotiques} (théorème~\ref{thm:nonasymp}) : $|s_n - s^*| \leq \Lambda^n |s_0 - s^*|$ pour tout $n$, avec $\Lambda$ explicite.
  \item \textbf{Double monotonie} (théorème~\ref{thm:monotonicity}) : $\lambda_{p,x}$ est croissant en $p$ et en $x$.
  \item \textbf{Extension au cas $x < 1$} (proposition~\ref{prop:x_small}) : la méthode converge avec un coefficient de poussière \emph{oscillatoire} $\lambda < 0$, avec $|\lambda| < 1$.
  \item \textbf{Point de départ optimal} (proposition~\ref{prop:optimal_start}) : $s_0 = h(1) = 1 - (x-1)/(xp)$ minimise le résidu initial.
  \item \textbf{Accélération de Steffensen} (théorème~\ref{thm:steffensen}) : le procédé $\Delta^2$ d'Aitken transforme la méthode linéaire de Pandrosion en une méthode à convergence quadratique sans dérivée, avec constante $K_S \ll K_{\text{Newton}}$.
  \item \textbf{Optimisation par changement d'échelle} (proposition~\ref{prop:scaling}) : en exploitant la monotonie de $\lambda_{p,x}$, un simple changement de variable réduit les problèmes à grand $x$ au cas $x' \approx 1$, divisant le coefficient de poussière par jusqu'à $39\times$ et effondrant le nombre d'itérations de centaines à quelques unités.
\end{itemize}

\medskip

La coexistence de \emph{cinq} profils de poussière distincts (Pandrosion, dichotomie, sécante, Newton, Steffensen--Pandrosion) pour la même cible $\alpha = \cbrt{2}$ illustre la thèse centrale : \textbf{la poussière est une propriété de la méthode, non du nombre.}

\medskip

\textbf{Perspectives.} Il serait naturel d'étudier des constructions de Pandrosion généralisées non linéaires (avec d'autres polynômes remplaçant $S_p$), d'explorer si des variantes de Steffensen d'ordre supérieur produisent des convergences encore plus rapides, ou d'analyser la structure topologique du bassin d'attraction pour $x < 1$ en dimension supérieure. La relation étroite entre l'itération de Pandrosion et l'arithmétique des racines suggère également des connexions avec les algorithmes de fractions continues qui méritent d'être explorées.

%=====================================================================
\begin{thebibliography}{9}

\bibitem{mactutor_pandrosion}
J.~J. O'Connor and E.~F. Robertson,
``Pandrosion of Alexandria,''
\textit{MacTutor History of Mathematics}, University of St~Andrews, 2020.
\url{https://mathshistory.st-andrews.ac.uk/Biographies/Pandrosion/}

\bibitem{mactutor_cube}
J.~J. O'Connor and E.~F. Robertson,
``Doubling the cube,''
\textit{MacTutor History of Mathematics}, University of St~Andrews, 2002.
\url{https://mathshistory.st-andrews.ac.uk/HistTopics/Doubling_the_cube/}

\bibitem{pappus}
Pappus d'Alexandrie,
\textit{Collectio}, Livre~III (circa 340~apr.~J.-C.).
Édition critique~: F.~Hultsch, \textit{Pappi Alexandrini Collectionis quae supersunt},
Weidmann, Berlin, 1876--1878.

\bibitem{wantzel}
P.~Wantzel,
``Recherches sur les moyens de reconnaître si un problème de Géométrie peut se résoudre avec la règle et le compas,''
\textit{Journal de Mathématiques Pures et Appliquées}, sér.~1, vol.~2, pp.~366--372, 1837.

\bibitem{stoer}
J.~Stoer and R.~Bulirsch,
\textit{Introduction to Numerical Analysis},
3rd~ed., Springer, New York, 2002.
Chap.~5~: itérations de point fixe, taux de convergence et accélération d'Aitken.

\bibitem{numrecipes}
W.~H. Press, S.~A. Teukolsky, W.~T. Vetterling, and B.~P. Flannery,
\textit{Numerical Recipes: The Art of Scientific Computing},
3rd~ed., Cambridge University Press, 2007.

\bibitem{steffensen}
J.~F. Steffensen,
``Remarks on iteration,''
\textit{Skandinavisk Aktuarietidskrift}, vol.~16, pp.~64--72, 1933.

\bibitem{aitken}
A.~C. Aitken,
``On Bernoulli's numerical solution of algebraic equations,''
\textit{Proceedings of the Royal Society of Edinburgh}, vol.~46, pp.~289--305, 1926.

\end{thebibliography}

\end{document}
